% 学术报告骨架（中文）。用 xelatex 编译：
% python scripts/compile_latex.py main.tex
\documentclass[12pt]{ctexrep}

\usepackage[margin=2.5cm]{geometry}
\usepackage{amsmath,amssymb}
\usepackage{graphicx}
\usepackage{booktabs}
\usepackage[hidelinks]{hyperref}

\title{学术报告标题}
\author{作者姓名}
\date{\today}

\begin{document}

% 封面
\begin{titlepage}
\centering
\vspace*{3cm}
{\LARGE\bfseries 学术报告标题\par}
\vspace{2cm}
{\large 作者姓名\par}
\vspace{1cm}
{\large 单位\par}
\vspace{1cm}
{\large \today\par}
\vfill
\end{titlepage}

% 摘要与关键词
\newpage
\noindent\textbf{摘要}\par
\vspace{0.5em}
摘要四句：背景、问题、内容、结论。
\par\vspace{1em}
\noindent\textbf{关键词：}关键词1、关键词2、关键词3

% 目录
\tableofcontents

\chapter{引言}
\label{chap:intro}
背景与目的。

\chapter{现状与数据}
\label{chap:data}
事实与数据，每个数字注明来源。

\chapter{分析}
\label{chap:analysis}
按问题分节，先给结论，后给依据。

\chapter{结论与建议}
\label{chap:conclusion}
结论。建议逐条列出，给出理由与预期效果。

\begin{thebibliography}{9}
\bibitem{example2026} 作者. 文献标题. 期刊名, 2026.
\end{thebibliography}

\appendix
\chapter{附录}
\label{app:detail}
补充材料。

\end{document}
